\documentclass[11pt]{letter}
\usepackage[margin=1in]{geometry}
\usepackage{parskip}

\signature{Hikaru Wakaura \\ Taiki Tanimae \\ QIRI (Quantum Integrated
Research Institute Inc.)}
\address{QIRI (Quantum Integrated Research Institute Inc.)\\
1--16--3 Akasaka, Minato-ku\\ Tokyo 107--0061, Japan\\
h.wakaura@deeptell.jp,\ t.tanimae@deeptell.jp}

\begin{document}
\begin{letter}{Editor-in-Charge \\ Physical Review A \\ American Physical
Society}

\opening{Dear Editor-in-Charge,}

We thank you for your decision and for taking the time to explain it. We
respectfully ask that you consider a substantially revised version of the
manuscript, ``Which quantum algorithms suit photonic continuous-variable
hardware? A comparative resource and noise-robustness study,'' attached
here, before our final redirection to another venue. We accept your
earlier judgement that the original submission did not contain a
sufficiently concrete physics result for PRA; in the revision we have
addressed that point directly with new physics content, summarised below.

\textbf{What is new since the previous submission.}

\textit{Worked physics example on three diatomic molecules
(new Sec.~VIII, new Fig.~4, new Table~II).} We compute the
vibrational eigenstructure of H$_2$, HF, and I$_2$ from their
Morse potentials in second-quantised form, using experimentally fixed
Huber--Herzberg spectroscopic constants (no free parameters). Control-free
CV quantum phase estimation, applied to the resulting Morse Hamiltonians,
recovers the populated vibrational eigenenergies of all three molecules
to within $0.55$~cm$^{-1}$, i.e.\ \emph{below the $1$~cm$^{-1}$
spectroscopic threshold standard in molecular physics}.  The recovery
spans a decade in vibrational quantum ($\omega_e=214$--$4401$~cm$^{-1}$)
and an order of magnitude in anharmonicity ($x_e=0.003$--$0.028$); the
weakly anharmonic case (I$_2$) reaches $\sim 10^{-2}$~cm$^{-1}$,
roughly two orders of magnitude inside the spectroscopic threshold.

\textit{Quantitative resource statement for a real physical observable.}
On the same H$_2$ vibrational spectrum, reaching $0.5$~cm$^{-1}$
accuracy with discrete-variable QPE requires roughly $1.5\times10^{10}$
T-gates and $1.6\times10^4$ controlled time evolutions; the
CV control-free protocol needs \emph{zero} controlled unitaries and
\emph{zero} non-Gaussian gates for the harmonic part. The advantage is
not incremental; it is structural and quantitative on a
\emph{physically observable spectrum}.

\textit{Validation and reproducibility.} The Hamiltonian construction
is exact in the truncated Fock basis (verified eigenvalue agreement
with the analytic Morse formula to machine precision at suitable
cutoff), the recovery is performed by reference-eigenstate
holographic phase retrieval, and every reported number is reproduced
by a single self-contained script bundled with the submission.

\textbf{Position relative to PRA scope.} Computing the vibrational
energy levels of real diatomic molecules to spectroscopic accuracy is, in
our view, a physical observable squarely within PRA's scope; the
algorithmic infrastructure that surrounds it is the means by which that
observable is computed. The three-molecule demonstration is not a single
toy example: it spans a wide swath of the molecular vibrational
parameter space. Beyond the methodology contribution of the original
manuscript, the new section provides a quantitative resource
prediction that experimentalists can directly test on near-term
GKP-encoded photonic processors.

We are aware that the threshold for PRA is high and that the editor's
discretion is final. We submit this revised version not to dispute the
earlier judgement on the original manuscript, but to ask whether the
addition of the H$_2$/HF/I$_2$ spectroscopic-accuracy calculation
changes the assessment. If it does not, we will redirect the work to
another venue without further appeal, and we are grateful for the
careful consideration given so far.

The manuscript has not been published and is not under consideration
elsewhere; the authors declare no competing interests. All scripts
and numerical results that reproduce every figure and table are
included with the submission.

\closing{Sincerely,}

\end{letter}
\end{document}
